%==============================================================================
\chapter{Resultados y análisis}
\label{cap:resultados}
%==============================================================================
% Objetivos específicos 1 y 4. Actividad A16.

La pregunta de investigación plantea determinar en qué medida un
\textit{schedule} basado en portafolios y guiado por aprendizaje automático
mejora la eficiencia y la robustez en el JSSP clásico. En este trabajo, esa
pregunta se responde para un portafolio concreto, formado por tres
configuraciones de Chuffed y Gecode (Chuffed, Gecode y Chuffed V2) con el
modelo con cotas, de modo que los resultados y las conclusiones se refieren a
ese portafolio y a las instancias estudiadas. Para responderla en el
orden de las cinco preguntas de la Sección~\ref{sec:preguntas}, se examinan a
continuación los resultados obtenidos sobre las 236 instancias con $lb < ub$
bajo un presupuesto de 300 segundos por instancia y validación
\textit{leave-one-out}, salvo que se indique otra cosa. Todas las tablas y
figuras se generan con el programa \texttt{src/09\_tablas\_figuras.py}.

%------------------------------------------------------------------------------
\section{Efecto de las cotas}
\label{sec:res_cotas}
%------------------------------------------------------------------------------

En la Tabla~\ref{tab:efecto_cotas} se contrasta el desempeño de Chuffed, Gecode
y CP-SAT entre el modelo original y el modelo con cotas.

\begin{table}[H]
  \centering
  \small
  \singlespacing
  \caption{Efecto de las cotas en el modelo (236 instancias, 300\,s).}
  \label{tab:efecto_cotas}
  \begin{tabular}{lrrrrrrrrr}
    \toprule
    \textit{Solver} & \multicolumn{2}{c}{Gap medio (\%)} & \multicolumn{3}{c}{Óptimos} & \multicolumn{2}{c}{$t$ mediano (s)} & Mejora & Empeora \\
    \cmidrule(lr){2-3}\cmidrule(lr){4-6}\cmidrule(lr){7-8}
     & Original & Cotas & Original & Cotas & Nuevos & Original & Cotas &  &  \\
    \midrule
    Chuffed & 648.3 & 16.8 & 28 & 36 & 8 & 60.1 & 2.8 & 207 & 1 \\
    Gecode & 25.6 & 21.7 & 5 & 6 & 1 & 2.6 & 0.4 & 172 & 51 \\
    CP-SAT (ref.) & 4.3 & 4.0 & 63 & 62 & 1 & 9.8 & 9.9 & 82 & 22 \\
    \bottomrule
  \end{tabular}
  \par\vspace{2pt}{\footnotesize $t$ mediano: tiempo hasta probar el óptimo en las instancias donde ambos modelos lo prueban. Mejora/Empeora: instancias con menor/mayor \textit{makespan} que con el modelo original.}
\end{table}


En Chuffed, la incorporación de las cotas modifica de manera sustancial el
desempeño. Su \textit{gap} medio pasa de 648.3\,\% a 16.8\,\%, puesto que con el
modelo original se quedaba en programas casi secuenciales en las instancias
grandes, mientras que con las cotas parte de soluciones al menos tan buenas como
la heurística. Prueba además la optimalidad en 8 instancias adicionales y, en
las instancias en que ambos modelos la prueban, el tiempo mediano baja de 60.1 a
2.8 segundos. Mejora en 207 de las 236 instancias y empeora solo en una.

En Gecode el efecto es más moderado. Su \textit{gap} medio baja de 25.6\,\% a
21.7\,\% y mejora en 172 instancias, pero empeora en 51. La razón obedece a la
segunda consecuencia de las cotas descrita en la
Sección~\ref{sec:modelo_cotas}, ya que Gecode no encuentra en 300 segundos un
programa mejor que el heurístico en 65 instancias, y en ellas su resultado es
la $ub$, a veces peor que la solución que encontraba con el modelo original.

En CP-SAT, que ya tenía un buen desempeño con el modelo original, la variación
es pequeña, con una reducción del \textit{gap} medio de 4.3\,\% a 4.0\,\%,
mejoras en 82 instancias, retrocesos en 22 y prácticamente los mismos óptimos
probados. CP-SAT calcula internamente sus propias cotas y soluciones iniciales,
de modo que la información adicional del modelo le aporta poco.

De este modo, la formulación con cotas sitúa a Chuffed como una configuración
competitiva dentro del portafolio y sirve de base para las evaluaciones
siguientes.

%------------------------------------------------------------------------------
\section{Desempeño de las configuraciones}
\label{sec:res_solvers}
%------------------------------------------------------------------------------

La Tabla~\ref{tab:individual} resume el desempeño de las tres configuraciones
del portafolio con un presupuesto de 300 segundos.

\begin{table}[H]
  \centering
  \small
  \singlespacing
  \caption{Desempeño de las configuraciones con el modelo con cotas (236 instancias, 300\,s).}
  \label{tab:individual}
  \begin{tabular}{lrrrrrr}
    \toprule
     & Chuffed & Gecode & Chuffed V2 & VBS & Heurística & CP-SAT (ref.) \\
    \midrule
    Óptimos probados & 36 & 6 & 18 & 36 & -- & 62 \\
    Gap medio (\%) & 16.76 & 21.68 & 15.08 & 12.83 & 26.96 & 3.97 \\
    PAR10 (s) & 2550 & 2924 & 2774 & 2550 & -- & 2224 \\
    Instancias donde gana & 100 & 3 & 100 & -- & -- & -- \\
    Sin solución & 12 & 65 & 24 & -- & -- & 0 \\
    \bottomrule
  \end{tabular}
  \par\vspace{2pt}{\footnotesize Se excluyen las 6 instancias con lb $=$ ub. Sin solución: el resultado es la ub. VBS: la mejor configuración del portafolio en cada instancia.}
\end{table}


Los resultados muestran un reparto claro de fortalezas entre las tres
configuraciones. Chuffed prueba la optimalidad en 36 instancias, el doble que
Chuffed V2, y encuentra solución en casi todas. Chuffed V2 tiene el menor
\textit{gap} medio (15.08\,\%) y constituye, por tanto, el SBS, al obtener la
mejor solución en 100 instancias, tantas como Chuffed, aunque prueba menos
óptimos. Gecode es la configuración más débil, ya que solo gana en 3 instancias
y no encuentra solución en 65. El VBS alcanza un \textit{gap} medio de
12.83\,\%, de modo que el margen disponible para el portafolio es de 2.24
puntos. La heurística sola tiene un \textit{gap} de 26.96\,\%.

\subsection{Complementariedad por tamaño y familia}

Al desagregar los resultados por número de operaciones
(Tabla~\ref{tab:por_tamano}) y cruzar el \textit{gap} de ambas versiones de
Chuffed (Figura~\ref{fig:dispersion}), se aprecia una complementariedad
estrechamente ligada al tamaño de la instancia.

\begin{table}[H]
  \centering
  \small
  \singlespacing
  \caption{Gap medio (\%) por tamaño de instancia.}
  \label{tab:por_tamano}
  \begin{tabular}{lrrrrrrrrr}
    \toprule
    Tamaño & Inst. & \multicolumn{5}{c}{Gap medio (\%)} & \multicolumn{3}{c}{Gana} \\
    \cmidrule(lr){3-7}\cmidrule(lr){8-10}
     &  & C & G & C2 & VBS & Margen & C & G & C2 \\
    \midrule
    6$\times$6 & 1 & 0.00 & 0.00 & 0.00 & 0.00 & 0.00 & 0 & 0 & 0 \\
    10$\times$5 & 5 & 0.00 & 4.39 & 1.91 & 0.00 & 0.00 & 0 & 0 & 0 \\
    15$\times$5 & 3 & 0.00 & 1.44 & 0.00 & 0.00 & 0.00 & 0 & 0 & 0 \\
    10$\times$10 & 18 & 0.07 & 9.04 & 3.15 & 0.07 & 0.00 & 9 & 0 & 0 \\
    20$\times$5 & 4 & 0.60 & 9.65 & 9.16 & 0.60 & 0.00 & 3 & 0 & 0 \\
    15$\times$10 & 5 & 0.21 & 15.39 & 2.85 & 0.21 & 0.00 & 4 & 0 & 0 \\
    20$\times$10 & 10 & 8.51 & 22.72 & 18.39 & 8.45 & 0.06 & 9 & 0 & 1 \\
    15$\times$15 & 15 & 2.10 & 16.08 & 6.01 & 2.10 & 0.00 & 15 & 0 & 0 \\
    30$\times$10 & 5 & 1.82 & 9.20 & 6.09 & 1.82 & 0.00 & 5 & 0 & 0 \\
    20$\times$15 & 28 & 12.29 & 24.23 & 16.63 & 11.86 & 0.43 & 22 & 0 & 6 \\
    20$\times$20 & 24 & 12.67 & 22.31 & 12.29 & 10.57 & 1.72 & 13 & 0 & 11 \\
    30$\times$15 & 20 & 24.91 & 28.57 & 21.26 & 19.50 & 1.76 & 8 & 0 & 12 \\
    50$\times$10 & 8 & 30.62 & 23.44 & 14.54 & 14.54 & 0.00 & 0 & 0 & 7 \\
    40$\times$15 & 10 & 32.06 & 29.46 & 23.02 & 22.89 & 0.12 & 2 & 0 & 8 \\
    30$\times$20 & 20 & 27.05 & 29.73 & 22.82 & 21.79 & 1.04 & 4 & 1 & 15 \\
    50$\times$15 & 20 & 25.18 & 22.81 & 17.87 & 17.64 & 0.23 & 3 & 1 & 14 \\
    40$\times$20 & 10 & 37.00 & 32.61 & 24.05 & 24.05 & 0.00 & 0 & 0 & 10 \\
    50$\times$20 & 20 & 29.27 & 27.05 & 21.99 & 21.94 & 0.05 & 3 & 1 & 15 \\
    100$\times$20 & 10 & 11.26 & 11.26 & 11.02 & 11.02 & 0.00 & 0 & 0 & 1 \\
    \bottomrule
  \end{tabular}
  \par\vspace{2pt}{\footnotesize C: Chuffed; G: Gecode; C2: Chuffed V2. Margen: gap de la mejor configuración individual del grupo menos gap del VBS.}
\end{table}


En las instancias de menos de 300 operaciones, Chuffed es mejor o igual que
las demás en casi todos los casos (Chuffed V2 gana solo en una), porque prueba
la optimalidad de la mayoría de las pequeñas y su búsqueda libre supera a la
estrategia de Chuffed V2. En $20 \times 15$ y $20 \times 20$ las victorias se
reparten, y a partir de $30 \times 15$ Chuffed V2 gana en la mayoría de las
instancias, con diferencias de hasta 16 puntos de \textit{gap} medio (en
$50 \times 10$). Gecode solo gana en tres instancias grandes de la familia
\texttt{ta}.

\begin{figure}[H]
  \centering
  \includegraphics{dispersion_configuraciones.pdf}
  \caption{\textit{Gap} de Chuffed frente al de Chuffed V2 en cada instancia,
  coloreado según la configuración que obtiene la mejor solución. Los puntos
  bajo la diagonal son instancias en que Chuffed V2 es mejor.}
  \label{fig:dispersion}
\end{figure}

Un comportamiento análogo aparece al agrupar las instancias por familia en la
Tabla~\ref{tab:por_familia}, donde Chuffed domina en las familias de instancias
pequeñas y medianas (\texttt{ft}, \texttt{la}, \texttt{abz}, \texttt{orb}), las
dos configuraciones de Chuffed se reparten las victorias en \texttt{swv},
\texttt{yn} y \texttt{ta}, y Chuffed V2 domina en \texttt{dmu}, con 65 de 80
victorias.

\begin{table}[H]
  \centering
  \small
  \singlespacing
  \caption{Resultados por familia de instancias.}
  \label{tab:por_familia}
  \begin{tabular}{lrrrrrrrrrrrrrr}
    \toprule
    Familia & Inst. & \multicolumn{3}{c}{Óptimos} & \multicolumn{4}{c}{Gap medio (\%)} & \multicolumn{3}{c}{Gana} & \multicolumn{3}{c}{Sin solución} \\
    \cmidrule(lr){3-5}\cmidrule(lr){6-9}\cmidrule(lr){10-12}\cmidrule(lr){13-15}
     &  & C & G & C2 & C & G & C2 & VBS & C & G & C2 & C & G & C2 \\
    \midrule
    ft & 3 & 2 & 1 & 1 & 0.80 & 12.72 & 7.77 & 0.80 & 2 & 0 & 0 & 0 & 0 & 0 \\
    la & 36 & 23 & 5 & 13 & 1.11 & 9.77 & 4.15 & 1.10 & 20 & 0 & 1 & 0 & 4 & 2 \\
    abz & 5 & 2 & 0 & 1 & 3.81 & 15.97 & 7.89 & 3.72 & 3 & 0 & 1 & 0 & 0 & 0 \\
    orb & 10 & 9 & 0 & 3 & 0.13 & 9.19 & 5.08 & 0.13 & 7 & 0 & 0 & 0 & 0 & 0 \\
    swv & 18 & 0 & 0 & 0 & 21.08 & 27.17 & 19.51 & 13.80 & 9 & 0 & 8 & 0 & 3 & 2 \\
    yn & 4 & 0 & 0 & 0 & 8.23 & 20.87 & 9.46 & 7.34 & 2 & 0 & 2 & 0 & 0 & 0 \\
    ta & 80 & 0 & 0 & 0 & 13.79 & 20.54 & 15.01 & 12.81 & 43 & 3 & 23 & 11 & 36 & 16 \\
    dmu & 80 & 0 & 0 & 0 & 29.70 & 29.25 & 21.32 & 20.81 & 14 & 0 & 65 & 1 & 22 & 4 \\
    \bottomrule
  \end{tabular}
  \par\vspace{2pt}{\footnotesize C: Chuffed; G: Gecode; C2: Chuffed V2.}
\end{table}


A partir de esta distribución se observa que acertar en la elección entre
Chuffed y Chuffed V2 basta para capturar casi todo el margen disponible, dado
que Gecode aporta muy poco al VBS. Sin embargo, la frontera entre ambas
configuraciones carece de un corte único, puesto que en los tamaños de
$20 \times 15$ a $30 \times 15$ las victorias están repartidas y las
características de la instancia, adicionales al tamaño, pueden inclinar la
decisión.

\subsection{Comportamiento en el tiempo}

La evolución del \textit{gap} medio de cada configuración según el tiempo
transcurrido puede seguirse en la Tabla~\ref{tab:anytime} y en la
Figura~\ref{fig:anytime}.

\begin{table}[H]
  \centering
  \small
  \singlespacing
  \caption{Gap medio (\%) y óptimos probados según el tiempo transcurrido.}
  \label{tab:anytime}
  \begin{tabular}{rrrrrrrr}
    \toprule
    Tiempo (s) & \multicolumn{4}{c}{Gap medio (\%)} & \multicolumn{3}{c}{Óptimos probados} \\
    \cmidrule(lr){2-5}\cmidrule(lr){6-8}
     & C & G & C2 & VBS & C & G & C2 \\
    \midrule
    10 & 22.19 & 22.25 & 16.82 & 16.02 & 22 & 4 & 9 \\
    30 & 20.53 & 21.99 & 16.11 & 15.04 & 26 & 5 & 13 \\
    60 & 19.36 & 21.94 & 15.76 & 14.46 & 28 & 5 & 16 \\
    120 & 18.21 & 21.79 & 15.42 & 13.77 & 30 & 6 & 16 \\
    300 & 16.76 & 21.68 & 15.08 & 12.83 & 36 & 6 & 18 \\
    \bottomrule
  \end{tabular}
  \par\vspace{2pt}{\footnotesize C: Chuffed; G: Gecode; C2: Chuffed V2.}
\end{table}


\begin{figure}[H]
  \centering
  \includegraphics{anytime.pdf}
  \caption{\textit{Gap} medio de cada configuración y del VBS en función del
  tiempo (escala logarítmica).}
  \label{fig:anytime}
\end{figure}

Chuffed V2 encuentra buenas soluciones muy pronto, pues a los 10 segundos su
\textit{gap} medio ya es de 16.8\,\%, cercano a su valor final, mientras que
Chuffed necesita los 300 segundos para llegar a 16.8\,\%. Gecode casi no
mejora después de los primeros segundos. Las dos configuraciones de Chuffed,
en cambio, continúan mejorando hasta el final del presupuesto, lo que anticipa
una pérdida de calidad potencial para cualquier estrategia que fraccione el
tiempo entre ellas.

%------------------------------------------------------------------------------
\section{\textit{Winner Takes All} frente al \textit{schedule} secuencial}
\label{sec:res_wta_schedule}
%------------------------------------------------------------------------------

La Tabla~\ref{tab:portafolios} y la Figura~\ref{fig:portafolios} comparan
todas las estrategias.

\subsection{Todas las estrategias superan al mejor \textit{solver} individual}

Todos los métodos de portafolio cierran entre el 76\,\% y el 87\,\% de la
brecha entre el SBS y el VBS, y la mejora frente al SBS es
estadísticamente significativa en todos los casos (por ejemplo, 1.96 puntos
para WTA $k$-NN, con un intervalo de confianza del 95\,\% de
$[1.48, 2.49]$, y 1.80 puntos para el \textit{schedule} con ML, con
$[1.36, 2.30]$). Asimismo, todos los métodos recuperan los 36 óptimos de
Chuffed o casi todos, que el SBS (Chuffed V2) pierde a la mitad. Incluso la
regla por tamaño, que usa una sola característica, cierra el 76\,\% de la
brecha.

\begin{table}[H]
  \centering
  \small
  \singlespacing
  \caption{Comparación de los \textit{solvers} sin cotas y con cotas y de los enfoques de portafolio (236 instancias, 300\,s por instancia, validación \textit{leave-one-out}).}
  \label{tab:portafolios}
  \begin{tabular}{lrrrrrr}
    \toprule
    Método & Gap (\%) & Óptimos & PAR10 (s) & Brecha cerrada (\%) & Mejor & Peor \\
    \midrule
    Chuffed (sin cotas) & 648.32 & 28 & 2652 & -- & 13 & 205 \\
    Gecode (sin cotas) & 25.57 & 5 & 2938 & -- & 7 & 225 \\
    VBS (sin cotas) & 24.53 & 28 & 2651 & -- & 17 & 201 \\
    \midrule
    Chuffed & 16.76 & 36 & 2550 & -- & 103 & 101 \\
    Gecode & 21.68 & 6 & 2924 & -- & 14 & 197 \\
    Chuffed V2 & 15.08 & 18 & 2774 & -- & 0 & 0 \\
    \midrule
    Regla por tamaño & 13.37 & 36 & 2550 & 76 & 68 & 8 \\
    WTA $k$-NN (mzn2feat) & 13.23 & 36 & 2550 & 82 & 84 & 18 \\
    WTA $k$-NN (propias) & 13.34 & 36 & 2550 & 77 & 83 & 17 \\
    WTA $k$-NN (ambas) & 13.12 & 36 & 2550 & 87 & 88 & 17 \\
    SUNNY (mzn2feat) & 13.20 & 36 & 2550 & 84 & 92 & 43 \\
    SUNNY (propias) & 13.28 & 34 & 2576 & 80 & 91 & 43 \\
    SUNNY (ambas) & 13.20 & 36 & 2552 & 84 & 92 & 40 \\
    \midrule
    \textit{Schedule} estático & 13.31 & 35 & 2567 & 79 & 97 & 70 \\
    \textit{Schedule} con ML (híbrido) & 13.28 & 35 & 2564 & 80 & 87 & 42 \\
    \midrule
    VBS & 12.83 & 36 & 2550 & 100 & 104 & 0 \\
    \bottomrule
  \end{tabular}
  \par\vspace{2pt}{\footnotesize Cada fila es un método. Gap: \textit{gap} medio frente a la mejor solución conocida. Óptimos: instancias en que se prueba la optimalidad. PAR10: tiempo medio hasta probar el óptimo, contando 3000\,s si no se prueba. En Gap y PAR10, un valor menor es mejor. Brecha cerrada: parte de la distancia entre el SBS (Chuffed V2) y el VBS que recupera el método, $100\cdot(\mathrm{gap}_{SBS}-\mathrm{gap})/(\mathrm{gap}_{SBS}-\mathrm{gap}_{VBS})$ (0 = SBS, 100 = VBS). Mejor/Peor: instancias con menor/mayor \textit{makespan} que el SBS. Sin cotas: Chuffed y Gecode con el modelo original y su VBS, es decir, el portafolio planteado al inicio del trabajo; Chuffed V2 no tiene fila sin cotas porque se definió como configuración del modelo con cotas. Las demás filas usan el modelo con cotas.}
\end{table}


Las filas sin cotas de la Tabla~\ref{tab:portafolios} corresponden al
portafolio planteado al inicio del trabajo, formado por Chuffed y Gecode con
el modelo original. Con ese modelo, el \textit{gap} medio de Chuffed era de
648.3\,\% y el de Gecode de 25.6\,\%, y aun eligiendo en cada instancia la
mejor de las dos, el VBS no bajaba del 24.5\,\%, casi el doble del 13\,\%
que alcanzan los enfoques de portafolio con el modelo con cotas. Además, el
margen entre el mejor \textit{solver} y el VBS era de apenas 1.0 puntos,
frente a 2.24 con cotas, de modo que el acotamiento del modelo es lo que le
da al portafolio espacio para mejorar. Chuffed V2 se incorporó después de
introducir las cotas, como una segunda estrategia de búsqueda sobre ese
modelo, por lo que no se evaluó con el modelo original.

\begin{figure}[H]
  \centering
  \includegraphics{portafolios_gap.pdf}
  \caption{\textit{Gap} medio de cada método. Las líneas punteadas marcan el
  SBS (gris) y el VBS (verde).}
  \label{fig:portafolios}
\end{figure}

\subsection{La selección supera en promedio al reparto del tiempo}

Con los datos de este trabajo, la mejor estrategia en promedio es
\textit{Winner Takes All} con las 85 características (13.12\,\%). El análisis
pareado frente a las demás estrategias, instancia por instancia
(Tabla~\ref{tab:comparaciones}), muestra que WTA obtiene un
\textit{makespan} igual o mejor que los \textit{schedules} en
la mayoría de las instancias en que difieren. La diferencia media frente a
SUNNY es pequeña y su intervalo de confianza incluye el cero, pero WTA es
mejor en 70 instancias y peor en 19. Frente al \textit{schedule} estático y
al \textit{schedule} con ML la ventaja es significativa, aunque de pocas
décimas.

\begin{table}[H]
  \centering
  \small
  \singlespacing
  \caption{Comparación pareada de WTA $k$-NN (ambas) con otros métodos en las 236 instancias.}
  \label{tab:comparaciones}
  \begin{tabular}{lrcr}
    \toprule
    Frente a & Diferencia & IC 95\,\% & WTA mejor / peor \\
    \midrule
    SUNNY (ambas) & $-0.08$ & $[-0.19, +0.04]$ & 70 / 19 \\
    \textit{Schedule} con ML (híbrido) & $-0.16$ & $[-0.28, -0.02]$ & 95 / 22 \\
    \textit{Schedule} estático & $-0.19$ & $[-0.33, -0.04]$ & 126 / 28 \\
    Regla por tamaño & $-0.25$ & $[-0.49, -0.03]$ & 26 / 15 \\
    \bottomrule
  \end{tabular}
  \par\vspace{2pt}{\footnotesize Diferencia: \textit{gap} medio de WTA menos el del otro método, en puntos porcentuales; un valor negativo favorece a WTA. IC 95\,\%: intervalo de confianza de esa diferencia, obtenido por \textit{bootstrap} con 10\,000 remuestreos de las instancias; si no incluye el cero, la diferencia es significativa. WTA mejor / peor: instancias en que WTA obtiene un \textit{makespan} menor / mayor que el otro método.}
\end{table}


Este comportamiento responde a dos factores observados en las secciones
previas. Primero, la selección sigue un patrón por tamaño muy marcado, lo que
permite a WTA obtener el mismo resultado que el VBS en 202 de las 236
instancias. Segundo, las dos configuraciones de Chuffed siguen mejorando hasta
los 300 segundos, de modo que reducir el tiempo asignado a la configuración
adecuada para destinarlo a otra disminuye la calidad alcanzada. Un
\textit{schedule} asume esa reducción en muchas instancias para protegerse de
un error de selección que, en este conjunto, ocurre pocas veces.

\subsection{El \textit{schedule} es más robusto}

El promedio constituye solo una dimensión del desempeño. La motivación de los
\textit{schedules} es reducir el riesgo de una mala selección, y la
Tabla~\ref{tab:robustez} muestra que en ese aspecto sí son mejores. Para cada
método se mide la pérdida frente al VBS, es decir, cuántos puntos de
\textit{gap} pierde en cada instancia respecto a la mejor configuración.

\begin{table}[H]
  \centering
  \small
  \singlespacing
  \caption{Pérdida frente al VBS (puntos de \textit{gap}) en las 236 instancias.}
  \label{tab:robustez}
  \begin{tabular}{lrrrrr}
    \toprule
    Método & Media & Percentil 90 & Máxima & $> 2$ puntos & $> 5$ puntos \\
    \midrule
    SBS (Chuffed V2) & 2.24 & 7.53 & 22.75 & 76 & 42 \\
    Regla por tamaño & 0.54 & 1.59 & 9.89 & 23 & 11 \\
    WTA $k$-NN (ambas) & 0.28 & 0.45 & 7.82 & 13 & 4 \\
    SUNNY (ambas) & 0.36 & 1.07 & 6.04 & 9 & 3 \\
    \textit{Schedule} con ML (híbrido) & 0.44 & 1.40 & 5.53 & 14 & 3 \\
    \textit{Schedule} estático & 0.47 & 1.28 & 4.31 & 9 & 0 \\
    \bottomrule
  \end{tabular}
\end{table}


WTA tiene la menor pérdida media, pero también la mayor pérdida máxima entre
los portafolios (7.82 puntos) y cuatro instancias en las que pierde más de 5
puntos, correspondientes a los casos en que elige la configuración equivocada.
Los \textit{schedules} cometen más errores pequeños, porque siempre ceden algo
de tiempo, pero limitan los grandes. El \textit{schedule} estático nunca pierde
más de 4.31 puntos, y SUNNY combina una pérdida media cercana a la de WTA con
menos instancias por encima de 2 puntos.

Respecto a la pregunta P3, la comparación refleja un compromiso entre el
desempeño promedio y el control del riesgo. Mientras que seleccionar una sola
configuración con WTA es ligeramente mejor en promedio, en el peor caso los
\textit{schedules} resultan más robustos, por lo cual la elección entre ambos
esquemas depende de si se prioriza el desempeño medio o acotar el efecto de una
mala decisión.

\subsection{Sensibilidad al número de vecinos}

Como se observa en la Tabla~\ref{tab:sensibilidad_k}, ni WTA ni SUNNY dependen
mucho de $k$, puesto que entre $k = 3$ y $k = 30$ el \textit{gap} medio varía
menos de 0.15 puntos. El valor $k = 15$ no se ajustó con estos resultados,
sino que se fijó de antemano como la raíz cuadrada del número de instancias
de entrenamiento, criterio que usa SUNNY-CP \citep{amadini2015sunnycp}, y
está entre los mejores para ambos métodos.

\begin{table}[H]
  \centering
  \small
  \singlespacing
  \caption{Sensibilidad al número de vecinos $k$ (características: ambas).}
  \label{tab:sensibilidad_k}
  \begin{tabular}{rrr}
    \toprule
    $k$ & Gap WTA (\%) & Gap SUNNY (\%) \\
    \midrule
    3 & 13.26 & 13.23 \\
    5 & 13.16 & 13.21 \\
    7 & 13.22 & 13.22 \\
    10 & 13.18 & 13.23 \\
    15 & 13.12 & 13.20 \\
    20 & 13.14 & 13.18 \\
    30 & 13.12 & 13.18 \\
    \bottomrule
  \end{tabular}
\end{table}


%------------------------------------------------------------------------------
\section{Aporte de las características y del aprendizaje automático}
\label{sec:res_prediccion}
%------------------------------------------------------------------------------

\subsection{Conjuntos de características}

En WTA, la unión de las características de \texttt{mzn2feat} y las propias
da el mejor resultado (13.12\,\%), seguida de \texttt{mzn2feat} sola
(13.23\,\%) y de las propias solas (13.34\,\%). Las diferencias son pequeñas y
afectan a pocas instancias, ya que frente a cada conjunto por separado el
resultado de la unión cambia en solo 11 instancias. En SUNNY, \texttt{mzn2feat}
y la unión empatan (13.20\,\%). Que las características genéricas y las propias
den resultados tan parecidos es coherente con la estructura de la
Sección~\ref{sec:res_solvers}, en la cual gran parte de la decisión depende del
tamaño, magnitud que ambos conjuntos describen bien.

Frente a la regla por tamaño, que usa solo el número de operaciones, WTA con
todas las características reduce el \textit{gap} en 0.25 puntos
(Tabla~\ref{tab:comparaciones}). Las características adicionales aportan,
sobre todo, en la zona de transición entre Chuffed y Chuffed V2.

\subsection{Calidad de los modelos de predicción}

En la Tabla~\ref{tab:ml_calidad} se detallan las métricas predictivas de los
tres modelos del \textit{schedule} con ML, medidas con sus predicciones fuera
de muestra.

\begin{table}[H]
  \centering
  \small
  \singlespacing
  \caption{Calidad de los modelos de aprendizaje automático (predicciones fuera de muestra, \textit{leave-one-out}).}
  \label{tab:ml_calidad}
  \begin{tabular}{llr}
    \toprule
    Modelo & Métrica & Valor \\
    \midrule
    Resolución (Chuffed) & AUC & 0.98 \\
    Resolución (Chuffed) & Exactitud (\%) & 95.3 \\
    Resolución (Gecode) & AUC & 0.98 \\
    Resolución (Gecode) & Exactitud (\%) & 97.0 \\
    Resolución (Chuffed V2) & AUC & 0.97 \\
    Resolución (Chuffed V2) & Exactitud (\%) & 94.1 \\
    Tiempo condicional (Chuffed) & Error absoluto medio (s) & 41.7 \\
    Tiempo condicional (Chuffed) & Error relativo mediano (\%) & 89.8 \\
    Mejor configuración & Exactitud (\%) & 76.7 \\
    Mejor configuración (Chuffed vs resto) & AUC & 0.85 \\
    Mejor configuración (Gecode vs resto) & AUC & 0.70 \\
    Mejor configuración (Chuffed V2 vs resto) & AUC & 0.86 \\
    \bottomrule
  \end{tabular}
\end{table}


El clasificador de resolución es muy preciso, con un área bajo la curva ROC
(AUC) de 0.97 a 0.98 y una exactitud de entre el 94\,\% y el 97\,\% para las
tres configuraciones, lo que indica que en este conjunto resulta sencillo
distinguir las instancias cuyos óptimos pueden probarse en 300 segundos. El
regresor de tiempo es el modelo más débil, con un error relativo mediano del
89.8\,\%, del mismo orden que el propio tiempo, razón por la cual el
\textit{schedule} asigna el doble del tiempo predicho. Su bajo desempeño
tiene varias causas posibles. Solo 60 de las 708 ejecuciones prueban el
óptimo y sirven para entrenarlo, casi todas en instancias pequeñas (100
operaciones de mediana, frente a 425 en el conjunto completo), y sus tiempos
van de 0.3 a 273 segundos. Además, el tiempo que tarda un \textit{solver} en
probar el óptimo depende en buena parte de su recorrido de búsqueda, que las
características estáticas de la instancia no describen, y las ejecuciones
que agotan los 300 segundos no indican cuánto habrían tardado. El clasificador de la
mejor configuración acierta en el 76.7\,\% de las instancias, con un AUC de
0.85 a 0.86 para las dos configuraciones de Chuffed y de 0.70 para Gecode, que
gana muy pocas veces.

\subsection{\textit{Schedule} con ML frente al \textit{schedule} estático}

La Tabla~\ref{tab:ml_vs_estatico} compara el \textit{schedule} con ML con el
\textit{schedule} estático, que reparte el tiempo igual en todas las
instancias.

\begin{table}[H]
  \centering
  \small
  \singlespacing
  \caption{\textit{Schedule} con ML frente al \textit{schedule} estático (diferencia ML $-$ estático).}
  \label{tab:ml_vs_estatico}
  \begin{tabular}{lr}
    \toprule
     & Valor \\
    \midrule
    Instancias con resultado distinto & 102 \\
    ML mejor / peor & 68 / 34 \\
    Diferencia media de gap (puntos) & -0.032 \\
    Intervalo de confianza del 95\,\% (\textit{bootstrap}) & [-0.145, +0.085] \\
    Mayor ganancia / mayor pérdida (puntos) & 4.24 / 4.98 \\
    Prueba de Wilcoxon ($p$) & 0.052 \\
    \bottomrule
  \end{tabular}
\end{table}


El \textit{schedule} con ML adapta el reparto a cada instancia (produce 181
\textit{schedules} distintos) y es mejor que el estático en 68 instancias y
peor en 34. Sin embargo, la diferencia media es de solo 0.03 puntos, su
intervalo de confianza incluye el cero y la prueba de Wilcoxon no es
significativa al 5\,\%. Los dos métodos son, por tanto, equivalentes en
promedio.

El \textit{schedule} con ML queda por detrás de WTA porque reparte el tiempo en
proporción a la probabilidad de que cada configuración sea la mejor, y esa
probabilidad rara vez es cercana a 1, de modo que incluso cuando el
clasificador acierta, cede parte del tiempo a las otras configuraciones. Se
probó también una versión del clasificador que pondera cada instancia de
entrenamiento por el costo de equivocarse, sin mejoras apreciables (13.26\,\%),
por lo que no se adoptó.

En relación con la pregunta P4, las características de la instancia aportan con
claridad, dado que todas las estrategias que las usan superan al SBS y a la
regla por tamaño. Aunque el aprendizaje automático produce modelos precisos, su
empleo para repartir el tiempo obtiene un desempeño medio equivalente al de un
reparto fijo bien elegido e inferior al de la selección por vecinos más
cercanos.

%------------------------------------------------------------------------------
\section{Comparación con el \textit{solver} de referencia}
\label{sec:res_referencia}
%------------------------------------------------------------------------------

CP-SAT se ejecutó con el modelo con cotas bajo el mismo protocolo que las
configuraciones del portafolio. Sus cifras globales ya se incluyeron como
referencia en la Tabla~\ref{tab:individual} (un \textit{gap} medio de
3.97\,\%, 62 óptimos probados y ninguna instancia sin solución), mientras que
el desglose frente al portafolio por número de operaciones se presenta en la
Tabla~\ref{tab:cpsat}.

\begin{table}[H]
  \centering
  \small
  \singlespacing
  \caption{CP-SAT frente al portafolio por número de operaciones (modelo con cotas, 300\,s).}
  \label{tab:cpsat}
  \begin{tabular}{lrrrrrrrr}
    \toprule
    Operaciones & Inst. & \multicolumn{4}{c}{Gap medio (\%)} & \multicolumn{2}{c}{Óptimos} & VBS mejor \\
    \cmidrule(lr){3-6}\cmidrule(lr){7-8}
     &  & CP-SAT & SBS & WTA $k$-NN & VBS & CP-SAT & VBS &  \\
    \midrule
    Hasta 150 & 36 & 0.00 & 3.25 & 0.13 & 0.13 & 34 & 32 & 0 \\
    151 a 300 & 58 & 1.49 & 13.28 & 8.20 & 7.88 & 18 & 4 & 0 \\
    301 a 600 & 82 & 4.13 & 18.58 & 17.91 & 17.37 & 7 & 0 & 0 \\
    Más de 600 & 60 & 8.52 & 19.13 & 19.13 & 19.04 & 3 & 0 & 1 \\
    \midrule
    Total & 236 & 3.97 & 15.08 & 13.12 & 12.83 & 62 & 36 & 1 \\
    \bottomrule
  \end{tabular}
  \par\vspace{2pt}{\footnotesize SBS: Chuffed V2. VBS mejor: instancias en que la mejor configuración del portafolio obtiene un \textit{makespan} menor que CP-SAT.}
\end{table}


CP-SAT supera ampliamente al portafolio. Su \textit{gap} medio es la tercera
parte del VBS (3.97\,\% frente a 12.83\,\%), prueba la optimalidad en 62
instancias frente a 36 y obtiene un \textit{makespan} igual o mejor que la
mejor configuración del portafolio en 235 de las 236 instancias, con la única
excepción de \texttt{dmu76}. La diferencia crece con el tamaño, ya que en las
instancias de hasta 150 operaciones todos los métodos están cerca del óptimo,
pero en las de más de 300 operaciones CP-SAT mantiene un \textit{gap} de 4 a
9\,\%, mientras que el portafolio se sitúa entre el 17 y el 19\,\%. Asimismo,
CP-SAT alcanza a los 10 segundos un \textit{gap} medio de 12.9\,\%, similar al
del portafolio a los 300 segundos.

Este resultado confirma lo observado con el modelo original
(Sección~\ref{sec:ejecucion_solvers}) y respalda la decisión de mantener a
CP-SAT fuera del portafolio, puesto que el VBS calculado junto con CP-SAT
coincide con este último en todas las instancias salvo una. La ventaja de
CP-SAT proviene de su arquitectura, que combina búsqueda con aprendizaje de
cláusulas, relajaciones lineales y búsqueda local sobre las soluciones
encontradas \citep{perron2023}. La comparación muestra, por tanto, que la
selección entre configuraciones de \textit{solvers} complementarios mejora de
forma consistente al mejor de ellos, aunque sin compensar la diferencia frente
a un \textit{solver} de nivel superior.

%------------------------------------------------------------------------------
\section{Discusión}
\label{sec:discusion}
%------------------------------------------------------------------------------

\subsection{Respuesta a las preguntas de evaluación}

\begin{enumerate}[label=\textbf{P\arabic*.}, leftmargin=*]
  \item \textbf{Efecto de las cotas.} Es el factor que más influye en el
    desempeño. Reduce el \textit{gap} de Chuffed de 648\,\% a 17\,\% y lo
    convierte en la base del portafolio, mientras que en Gecode la mejora es
    moderada.
  \item \textbf{Complementariedad.} Las configuraciones son complementarias y
    muestran un patrón bien definido, con ventaja de Chuffed en las instancias
    pequeñas y medianas, de Chuffed V2 en las grandes y un aporte marginal de
    Gecode. El margen entre el SBS y el VBS es de 2.24 puntos.
  \item \textbf{WTA frente a \textit{schedule}.} En promedio, con estas
    configuraciones, la selección de una configuración con WTA es la mejor estrategia (87\,\% de la brecha
    cerrada), en tanto que los \textit{schedules} quedan algo por detrás en el
    promedio pero resultan más robustos frente a errores de selección.
  \item \textbf{Características y aprendizaje automático.} Las características
    permiten cerrar entre el 77\,\% y el 87\,\% de la brecha. Los modelos de
    aprendizaje son precisos, pero el \textit{schedule} construido con ellos
    es equivalente al estático.
  \item \textbf{Estado del arte.} CP-SAT, usado como referencia externa,
    supera a este portafolio en casi todas las instancias (3.97\,\% de
    \textit{gap} frente a 13.12\,\%). El portafolio mejora a sus miembros,
    pero no alcanza a CP-SAT.
\end{enumerate}

\subsection{Amenazas a la validez}

\begin{itemize}
  \item El margen disponible entre el SBS y el VBS es de 2.24 puntos de
    \textit{gap}, y las diferencias entre los métodos de portafolio son de
    décimas, razón por la cual las comparaciones se apoyan en pruebas pareadas
    y no solo en promedios.
  \item Todos los resultados se obtienen sobre instancias clásicas de la
    literatura (236 en la evaluación del portafolio) y, aunque la validación \textit{leave-one-out} evita evaluar
    con instancias de entrenamiento, las conclusiones podrían variar en
    instancias de otra distribución o estructura.
  \item Los candidatos del \textit{schedule} estático y la regla de reparto del
    \textit{schedule} con ML se definieron tras analizar la campaña, por lo que
    aunque la elección final se calcula dentro de cada partición de
    entrenamiento, el conjunto de alternativas estuvo guiado por los datos.
  \item Las mejores soluciones conocidas (BKS) corresponden a los resultados
    recopilados hasta 2020 \citep{weise2019jssp}, por lo cual en las instancias
    abiertas el \textit{gap} es una cota inferior de la distancia real al
    óptimo, condición que afecta por igual a todos los métodos.
  \item Las configuraciones dentro de un \textit{schedule} se ejecutan sin
    compartir información, de modo que un portafolio que transmitiera la mejor
    solución de una configuración a la siguiente como nueva cota superior
    podría obtener mejores resultados que los reportados.
  \item Cada configuración se ejecutó una sola vez por instancia y, aunque los
    \textit{solvers} son deterministas con parámetros fijos, los tiempos de
    pared presentan pequeñas fluctuaciones que pueden alterar la última
    solución alcanzada cerca del límite.
\end{itemize}
